The current implementation supports two Stage 1 training modes:
\begin{enumerate}[leftmargin=1.25em]
  \item \textbf{Baseline.} Train only on answerable examples so the model is
  optimized for standard extractive QA behavior.
  \item \textbf{Abstain-aware.} Train on both answerable and unanswerable
  examples so that unsupported cases are represented during supervision.
\end{enumerate}

Both modes use the same underlying extractive QA model family and share the
same preprocessing and evaluation pipeline.
At inference time, the system converts start and end logits into either an
answer span or an abstention decision.
For the abstain-aware configuration, a validation-selected null-score threshold
is used to control the answer-versus-abstain trade-off.

The purpose of Stage 1 is not to introduce a complicated control mechanism.
Instead, it establishes whether explicit abstention training alone yields a
better supported-answering profile than an answer-only baseline under a tightly
controlled setting.

The broader roadmap already anticipates later stages for:
\begin{itemize}[leftmargin=1.25em]
  \item evidence support verification,
  \item confidence calibration,
  \item unsupported-confidence control,
  \item retrieval-grounded QA, and
  \item adaptive constraint balancing.
\end{itemize}

Those stages are part of the research direction, but only the Stage 1 path is
implemented end-to-end in the current codebase.
